
\documentclass[11pt]{article}

\usepackage[a4paper,margin=1in]{geometry}
\usepackage{amsmath,amssymb,mathtools}
\usepackage{hyperref}
\usepackage{xcolor}
\usepackage{listings}

\hypersetup{
  colorlinks=true,
  linkcolor=blue,
  urlcolor=blue,
  citecolor=blue
}

\lstset{
  basicstyle=\ttfamily\small,
  breaklines=true,
  columns=fullflexible,
  frame=single,
  backgroundcolor=\color{gray!5}
}

\title{\texttt{xyz2nlm}: Constructive Cartesian-to-Spherical\\
3D Harmonic-Oscillator Basis Transformation}
\author{}
\date{}

\begin{document}
\maketitle

\section{Purpose}

The package \texttt{xyz2nlm} constructs the spherical harmonic-oscillator basis
\[
    |N,L,M\rangle
\]
from the Cartesian harmonic-oscillator basis
\[
    |n_x,n_y,n_z\rangle .
\]
The shell quantum number is
\[
    N=n_x+n_y+n_z .
\]
In the code this same shell quantum number is stored as metadata named \texttt{n}.
Thus, throughout this document,
\[
    \texttt{state.n}=N .
\]

The algorithm is constructive. It does not diagonalize angular-momentum matrices.
Instead it performs two steps:

\begin{enumerate}
    \item construct the highest-weight states
    \[
        |N,L,L\rangle
    \]
    directly in the Cartesian basis;
    \item obtain the remaining states
    \[
        |N,L,M\rangle,\qquad M<L,
    \]
    by repeated application of the angular-momentum lowering operator
    \[
        \hat L_-=\hat L_x-i\hat L_y .
    \]
\end{enumerate}

The important implementation point is that all states live inside a fixed shell
\(N\), whose Cartesian basis dimension is
\[
    d_N=\frac{(N+1)(N+2)}{2}.
\]

\section{Cartesian basis indexing}

For fixed \(N\), define the set of Cartesian occupation configurations
\[
    \mathcal C_N
    =
    \left\{
        (n_x,n_y,n_z)\in\mathbb N_0^3:
        n_x+n_y+n_z=N
    \right\}.
\]
The number of such states is \(d_N=(N+1)(N+2)/2\).

The implementation uses a bosonic permanent, or stars-and-bars, indexing map
\[
    J_N:\mathcal C_N\to \{0,1,\dots,d_N-1\}.
\]
This is the index convention used by the \texttt{mlx\_create\_c} package,
following the bosonic occupation indexing described in Ref.~\cite{arxiv09102577}.
The inverse operation is used conceptually as
\[
    J_N^{-1}(j)=(n_x,n_y,n_z).
\]

A vector in the fixed-\(N\) Cartesian basis is therefore represented as
\[
    |\psi_N\rangle
    =
    \sum_{(n_x,n_y,n_z)\in\mathcal C_N}
    \psi_{n_xn_yn_z}\,
    |n_x,n_y,n_z\rangle ,
\]
or equivalently as a dense array
\[
    \psi[J_N(n_x,n_y,n_z)]
    =
    \psi_{n_xn_yn_z}.
\]

The code generally stores only nonzero components. The compressed
representation is
\[
    \bigl(\texttt{non\_zero\_ind},\texttt{non\_zero\_val}\bigr),
\]
where

\begin{itemize}
    \item \texttt{non\_zero\_ind} is an integer \texttt{np.ndarray} containing
    the nonzero indices \(J_N(n_x,n_y,n_z)\);
    \item \texttt{non\_zero\_val} is a complex \texttt{np.ndarray} containing
    the corresponding amplitudes.
\end{itemize}

Thus the compressed state means
\[
    |\psi_N\rangle
    =
    \sum_{r}
    \texttt{non\_zero\_val}[r]\,
    \left|
        J_N^{-1}\bigl(\texttt{non\_zero\_ind}[r]\bigr)
    \right\rangle .
\]

\section{State container classes}

The relevant state containers are defined in
\[
    \texttt{xyz2nlm.custom\_dataclasses}.
\]

\subsection{\texttt{TransientState}}

The class
\[
    \texttt{xyz2nlm.custom\_dataclasses.TransientState}
\]
is used for intermediate calculations. It stores

\[
    \texttt{n},\qquad
    \texttt{non\_zero\_ind},\qquad
    \texttt{non\_zero\_val}.
\]

The shell metadata \texttt{n} is required because the same integer index has
meaning only after specifying the shell \(N=\texttt{n}\). The dense
representation is recovered with
\[
    \texttt{state.as\_ndarray()}.
\]
This method creates a dense array of length
\[
    d_N=\frac{(N+1)(N+2)}{2}
\]
and fills
\[
    \texttt{out[state.non\_zero\_ind]}=
    \texttt{state.non\_zero\_val}.
\]

The class implements two algebraic operations.

\paragraph{Scaling.}
For \(c\in\mathbb C\),
\[
    c|\psi\rangle
\]
is represented by leaving \texttt{non\_zero\_ind} unchanged and replacing
\[
    \texttt{non\_zero\_val}
    \mapsto
    c\,\texttt{non\_zero\_val}.
\]

\paragraph{Addition.}
For two states in the same shell \(N\),
\[
    |\chi\rangle=|\psi\rangle+|\phi\rangle
\]
is implemented by merging the two index sets and summing amplitudes on
indices that occur in both inputs. In formula form,
\[
    \chi_j=\psi_j+\phi_j,
\]
but the implementation performs this operation directly on the compressed
index/value arrays.

\subsection{\texttt{SphericalHOBasisState}}

The class
\[
    \texttt{xyz2nlm.custom\_dataclasses.SphericalHOBasisState}
\]
stores a finalized spherical harmonic-oscillator state. It carries the same
compressed vector data as \texttt{TransientState}, but also records
\[
    \texttt{n}=N,\qquad
    \texttt{l}=L,\qquad
    \texttt{m}=M .
\]
It is not intended as the main arithmetic workhorse. Its role is to associate
a compressed Cartesian vector with the spherical-basis label
\[
    |N,L,M\rangle .
\]
It also supports extraction to a dense \texttt{np.ndarray} through
\[
    \texttt{state.as\_ndarray()}.
\]

\section{Highest-weight seed states}

\subsection{Definition}

The seed states are the highest-weight states
\[
    |N,L,L\rangle .
\]
They are constructed directly in
\[
    \texttt{xyz2nlm.direct\_seed\_state\_calculation.create\_seed\_state}.
\]

The construction is based on the operator identity
\[
    |N,L,L\rangle
    \propto
    \bigl(R_{1,1,1}\bigr)^L
    \bigl(S^\dagger\bigr)^q
    |0,0,0\rangle ,
    \qquad
    q=\frac{N-L}{2}.
\]
The two creation-only operators are
\[
    R_{1,1,1}
    =
    -\frac{\ell}{2}
    \left(a_x^\dagger+i a_y^\dagger\right),
\]
and
\[
    S^\dagger
    =
    a_x^{\dagger 2}
    +
    a_y^{\dagger 2}
    +
    a_z^{\dagger 2}.
\]
Because both are polynomials only in creation operators, all factors commute.
Therefore the expansion into Cartesian number states can be computed
combinatorially.

The implementation does not literally apply these operators. It evaluates
the closed-form Cartesian coefficients
\[
    |N,L,L\rangle
    =
    \sum_{n_x+n_y+n_z=N}
    C^{(N,L)}_{n_x,n_y,n_z}
    |n_x,n_y,n_z\rangle .
\]

\subsection{Admissible Cartesian configurations}

A Cartesian state contributes to \(|N,L,L\rangle\) only if

\[
    n_x+n_y+n_z=N,
\]
\[
    n_z\ \text{is even},
\]
and
\[
    0\le n_z\le N-L .
\]
The parity condition
\[
    N-L\in 2\mathbb Z
\]
is also required. Equivalently,
\[
    q=\frac{N-L}{2}\in\mathbb N_0 .
\]

For a valid configuration define
\[
    s=q-\frac{n_z}{2}
    =
    \frac{N-L-n_z}{2}
    =
    \frac{n_x+n_y-L}{2}.
\]

\subsection{The integer polynomial coefficient \(S_{L,s}(n)\)}

The nontrivial cancellation in the seed-state coefficient is isolated in the
integer coefficient
\[
    S_{L,s}(n).
\]
It is defined by the generating function
\[
    (1+x)^L(1-x^2)^s
    =
    \sum_{n=0}^{L+2s}
    S_{L,s}(n)x^n .
\]
Equivalently,
\[
    S_{L,s}(n)
    =
    \sum_{\alpha}
    (-1)^\alpha
    \binom{s}{\alpha}
    \binom{L}{n-2\alpha},
\]
where the binomial coefficient is interpreted as zero when its lower
argument is outside the allowed range.

The implementation computes \(S_{L,s}(n)\) with Python's arbitrary-precision
integers in
\[
    \texttt{xyz2nlm.direct\_seed\_state\_calculation.Sls}.
\]
It uses the recurrence
\[
    \boxed{
    S_{L,s}(n)
    =
    \frac{L}{n}S_{L,s}(n-1)
    -
    \frac{L+2s-n+2}{n}S_{L,s}(n-2)
    }
\]
for
\[
    2\le n\le L+2s,
\]
with initial values
\[
    S_{L,s}(0)=1,
    \qquad
    S_{L,s}(1)=L.
\]
For the special case \(L=s=0\), only \(S_{0,0}(0)=1\) is used.

This recurrence follows from
\[
    G_{L,s}(x)=(1+x)^L(1-x^2)^s
\]
and
\[
    (1-x^2)G'_{L,s}(x)
    =
    \left[L-(L+2s)x\right]G_{L,s}(x).
\]
Equating coefficients gives
\[
    nS_{L,s}(n)
    =
    L S_{L,s}(n-1)
    -
    (L+2s-n+2)S_{L,s}(n-2).
\]

\subsection{Closed-form seed coefficient}

The coefficient of
\[
    |n_x,n_y,n_z\rangle
\]
in \(|N,L,L\rangle\) is
\[
    \boxed{
    C^{(N,L)}_{n_x,n_y,n_z}
    =
    (-1)^L i^{L-n_x}
    S_{L,s}(n_x)
    \binom{q}{s}
    \sqrt{
    2^{-L}
    \frac{
        (2L+1)!\,(L+q)!\,n_x!\,n_y!\,n_z!
    }{
        (L!)^2\,q!\,(N+L+1)!
    }}
    }
\]
where
\[
    q=\frac{N-L}{2},
    \qquad
    s=q-\frac{n_z}{2}.
\]

The code uses the equivalent phase convention
\[
    (-1)^L i^{L-n_x}
    =
    i^{3L-n_x}.
\]
Therefore the coefficient is evaluated in the form
\[
    C^{(N,L)}_{n_x,n_y,n_z}
    =
    \operatorname{sgn}\!\left(S_{L,s}(n_x)\right)
    i^{3L-n_x}
    e^\Phi ,
\]
where
\[
\begin{aligned}
    \Phi
    ={}&
    \log\left|S_{L,s}(n_x)\right|
    -\frac{L}{2}\log 2
    \\
    &+
    \log\Gamma(q+1)
    -\log\Gamma(s+1)
    -\log\Gamma(q-s+1)
    \\
    &+
    \frac{1}{2}
    \Big[
        \log\Gamma(2L+2)
        +\log\Gamma(L+q+1)
        +\log\Gamma(n_x+1)
        +\log\Gamma(n_y+1)
        +\log\Gamma(n_z+1)
    \\
    &\hspace{3.5cm}
        -2\log\Gamma(L+1)
        -\log\Gamma(q+1)
        -\log\Gamma(N+L+2)
    \Big].
\end{aligned}
\]
This is implemented in
\[
    \texttt{xyz2nlm.direct\_seed\_state\_calculation.coefficient\_seed\_state\_expansion}.
\]
The logarithmic form is used to avoid overflow in factorials while still
returning an ordinary double-precision complex coefficient.

The function
\[
    \texttt{xyz2nlm.direct\_seed\_state\_calculation.create\_seed\_state}
\]
does the full job:
\begin{enumerate}
    \item generate all valid Cartesian configurations;
    \item compute \(S_{L,s}(n_x)\);
    \item compute the coefficient \(C^{(N,L)}_{n_x,n_y,n_z}\);
    \item map configurations to indices \(J_N(n_x,n_y,n_z)\);
    \item package the result as a
    \texttt{xyz2nlm.custom\_dataclasses.SphericalHOBasisState}.
\end{enumerate}

\section{Angular-momentum operators in the Cartesian basis}

The angular-momentum operators are implemented in
\[
    \texttt{xyz2nlm.create\_angular\_momentum\_matrices.CreateAngularMomentumMatrices}.
\]
The implementation uses units with
\[
    \hbar=1.
\]

The Cartesian formulas used are
\[
    L_x
    =
    -i\,a_y^\dagger a_z
    +
    i\,a_z^\dagger a_y,
\]
\[
    L_y
    =
    -i\,a_z^\dagger a_x
    +
    i\,a_x^\dagger a_z,
\]
and
\[
    L_z
    =
    -i\,a_x^\dagger a_y
    +
    i\,a_y^\dagger a_x.
\]

Therefore
\[
\begin{aligned}
    L_+
    &=
    L_x+iL_y \\
    &=
    -i\,a_y^\dagger a_z
    +
    i\,a_z^\dagger a_y
    +
    a_z^\dagger a_x
    -
    a_x^\dagger a_z,
\end{aligned}
\]
and
\[
\begin{aligned}
    L_-
    &=
    L_x-iL_y \\
    &=
    -i\,a_y^\dagger a_z
    +
    i\,a_z^\dagger a_y
    -
    a_z^\dagger a_x
    +
    a_x^\dagger a_z.
\end{aligned}
\]

The method
\[
    \texttt{get\_ladder\_operators()}
\]
returns
\[
    (L_+,L_-),
\]
with \(L_-\) as the second output.

\section{Sparse implementation of \(a_i^\dagger a_j\)}

The angular-momentum operators are sums of terms of the form
\[
    c_{ij}\,a_i^\dagger a_j,
    \qquad i,j\in\{x,y,z\}.
\]
The code constructs these terms without explicitly building dense matrices.
It uses the following identity.

Let
\[
    \mathbf k=(k_x,k_y,k_z),
    \qquad
    k_x+k_y+k_z=N-1.
\]
Then
\[
    a_i^\dagger a_j
    |\mathbf k+\mathbf e_j\rangle
    =
    \sqrt{(k_i+1)(k_j+1)}
    |\mathbf k+\mathbf e_i\rangle .
\]
Thus a term \(c_{ij}a_i^\dagger a_j\) contributes
\[
    H_{J_N(\mathbf k+\mathbf e_i),\,J_N(\mathbf k+\mathbf e_j)}
    \mathrel{+}=
    c_{ij}\sqrt{(k_i+1)(k_j+1)}.
\]

The core mapping procedure is:

\begin{lstlisting}[language=Python]
def _calculate_mapping_matrices(self):
    self.Ns1 = mlx_create_c.create_bos_ns(self.N-1, 3)
    self.mapmat = mlx_create_c.create_mapmat_bos(self.Ns1)
    self.basis_size = self.mapmat[-1,-1] + 1

def _calculate_linear_operator_elements(self, i, j):
    mapmat = self.mapmat.view()
    Ns1 = self.Ns1.view()

    lhs_map = mapmat[:, i]
    rhs_map = mapmat[:, j]
    coefficients = np.sqrt((Ns1[:, i] + 1.) * (Ns1[:, j] + 1.))

    return lhs_map, rhs_map, coefficients
\end{lstlisting}

Here:

\begin{itemize}
    \item \texttt{Ns1} stores all \((N-1)\)-shell configurations
    \(\mathbf k\);
    \item \texttt{mapmat[:, i]} stores
    \[
        J_N(\mathbf k+\mathbf e_i)
    \]
    for each \(\mathbf k\);
    \item \texttt{mapmat[:, j]} stores
    \[
        J_N(\mathbf k+\mathbf e_j);
    \]
    \item the coefficient array stores
    \[
        \sqrt{(k_i+1)(k_j+1)}.
    \]
\end{itemize}

The resulting operators are packaged as
\[
    \texttt{scipy.sparse.linalg.LinearOperator}
\]
instances with both \texttt{matvec} and \texttt{matmat} implemented. For an
input vector \(\psi\), a term \(c_{ij}a_i^\dagger a_j\) acts as
\[
    (H\psi)_{J_N(\mathbf k+\mathbf e_i)}
    \mathrel{+}=
    c_{ij}\sqrt{(k_i+1)(k_j+1)}
    \psi_{J_N(\mathbf k+\mathbf e_j)}.
\]
The same indexing rule is applied columnwise in \texttt{matmat}.

\section{Generating \(M<L\) states}

After a seed state \(|N,L,L\rangle\) is available, the remaining states in the
same multiplet are generated by repeated application of \(L_-\).

The angular-momentum normalization convention is
\[
    L_-|N,L,M\rangle
    =
    \sqrt{(L+M)(L-M+1)}
    |N,L,M-1\rangle .
\]
Therefore
\[
    |N,L,M-1\rangle
    =
    \frac{
        L_-|N,L,M\rangle
    }{
        \sqrt{(L+M)(L-M+1)}
    }.
\]

This is the normalization applied in
\[
    \texttt{BlockLmApplication.normalize\_states\_after\_lm\_application}.
\]
If a matrix contains several current states as columns, the implementation
uses \texttt{matmat} to apply \(L_-\) to all columns simultaneously, then
divides each column by its own factor
\[
    \sqrt{(L+M)(L-M+1)}.
\]

\section{Block orchestration}

The full basis generation is orchestrated by
\[
    \texttt{xyz2nlm.create\_spherical\_basis.SphericalHOBasis}.
\]
The user initializes
\[
    \texttt{SphericalHOBasis(Nmax)}
\]
and calls
\[
    \texttt{calculate\_states()}.
\]

The implementation allocates shell blocks for
\[
    0\le N\le N_{\max}.
\]
For each fixed \(N\), the allowed angular momenta are
\[
    L=N,N-2,N-4,\dots,
\]
down to \(0\) for even \(N\) and down to \(1\) for odd \(N\).
Equivalently,
\[
    0\le L\le N,
    \qquad
    N-L\in 2\mathbb Z .
\]

The workflow is:

\begin{enumerate}
    \item Create all highest-weight seed states
    \[
        |N,L,L\rangle
    \]
    using
    \[
        \texttt{create\_seed\_state(N,L)}.
    \]
    This step is done concurrently when possible.

    \item Once all seed states for a fixed shell \(N\) are available, stack
    them as columns of a matrix
    \[
        B_N
        =
        \begin{bmatrix}
        |N,L_1,L_1\rangle &
        |N,L_2,L_2\rangle &
        \cdots
        \end{bmatrix}.
    \]

    \item Obtain \(L_-\) for the same shell with
    \[
        \texttt{CreateAngularMomentumMatrices(N).get\_ladder\_operators()}.
    \]

    \item Repeatedly apply \(L_-\) to the active block of columns, normalize
    each column by
    \[
        \sqrt{(L+M)(L-M+1)},
    \]
    and store the resulting states
    \[
        |N,L,M-1\rangle.
    \]

    \item Drop a multiplet from the active block once its lowest state
    \(M=-L\) has been produced.
\end{enumerate}

This blockwise construction avoids applying \(L_-\) separately to every state.
For each shell \(N\), many states are updated simultaneously through the
\texttt{LinearOperator.matmat} implementation.

\section{Internal consistency checks}

The mathematically expected checks are:

\paragraph{Normalization.}
For every produced state,
\[
    \langle N,L,M|N,L,M\rangle=1.
\]

\paragraph{\(L_z\) eigenvalue.}
\[
    L_z|N,L,M\rangle=M|N,L,M\rangle .
\]

\paragraph{Highest-weight condition.}
For seed states,
\[
    L_+|N,L,L\rangle=0 .
\]

\paragraph{Casimir eigenvalue.}
\[
    L^2|N,L,M\rangle=L(L+1)|N,L,M\rangle .
\]
In code,
\[
    L^2=L_x^2+L_y^2+L_z^2.
\]
For numerical validation it is usually more stable to first check
\(L_+\) on highest-weight states and \(L_z\) on all states, because those
checks involve less cancellation than a direct Cartesian-basis application of
\(L^2\).

\section{Implementation summary for maintainers and LLMs}

The code should be understood as follows:

\begin{itemize}
    \item The package constructs a change-of-basis matrix from Cartesian
    harmonic-oscillator states to spherical harmonic-oscillator states.

    \item The shell quantum number \(N\) is called \texttt{n} in the code.
    A fixed-\(N\) Cartesian sector has dimension
    \[
        (N+1)(N+2)/2.
    \]

    \item A state is stored sparsely by \texttt{non\_zero\_ind} and
    \texttt{non\_zero\_val}. The indices are occupation-state indices
    \(J_N(n_x,n_y,n_z)\).

    \item \texttt{TransientState} is for arithmetic: scaling and addition.

    \item \texttt{SphericalHOBasisState} is for finalized basis states and
    carries the metadata \((N,L,M)\).

    \item Seed states \(|N,L,L\rangle\) are computed directly from a closed
    formula for
    \[
        C^{(N,L)}_{n_x,n_y,n_z}.
    \]
    No numerical diagonalization is involved.

    \item The integer factor \(S_{L,s}(n_x)\) is computed exactly using
    Python integers and the recurrence
    \[
        S_{L,s}(n)
        =
        \frac{L}{n}S_{L,s}(n-1)
        -
        \frac{L+2s-n+2}{n}S_{L,s}(n-2).
    \]

    \item The positive factorial prefactor in the coefficient is evaluated
    with \texttt{scipy.special.gammaln}.

    \item The states with \(M<L\) are obtained by applying the sparse
    \texttt{LinearOperator} for \(L_-\) and normalizing by
    \[
        \sqrt{(L+M)(L-M+1)}.
    \]

    \item The angular-momentum matrices are not formed densely. Each term
    \(a_i^\dagger a_j\) is implemented by mapping
    \[
        |\mathbf k+\mathbf e_j\rangle
        \mapsto
        |\mathbf k+\mathbf e_i\rangle
    \]
    for all \((N-1)\)-particle configurations \(\mathbf k\).
\end{itemize}

\begin{thebibliography}{9}

\bibitem{arxiv09102577}
A. I. Streltsov, O. E. Alon, and L. S. Cederbaum,
``General mapping for bosonic and fermionic operators in Fock space'',
Phys. Rev. A \textbf{81}, 022124 (2010).
\url{https://link.aps.org/doi/10.1103/PhysRevA.81.022124}

\bibitem{xyz2nlm}
G. M. Koutentakis,
\texttt{xyz2nlm}: Create the spherical harmonic oscillator basis from the
Cartesian one.
\url{https://github.com/gkoutentakis/xyz2nlm}

\bibitem{mlxcreatec}
G. M. Koutentakis,
\texttt{mlx\_create\_c}: Port of the MATLAB create library in C for use in
other programming languages.
\url{https://github.com/gkoutentakis/mlx_create_c}

\end{thebibliography}

\end{document}
